\clearpage
\section{关键注意事项}

\subsection{章节换页}

每一章正文结束后，建议在 \verb|\section| 之前使用 \verb|\clearpage| 命令强制换页，以确保图表等浮动体不会跨章出现。

\subsection{标题换行命令}

本模板定义了专用的标题换行命令 \verb|\titlebreak|\index{huanhangxingwei@换行行为}，用于在不同排版场景下独立控制标题的换行行为。\verb|\titlebreak| 与 \verb|\\| 的区别如表~\ref{tab:titlebreak} 所示。

\begin{table}[H]
    \centering
    \bicaption{标题换行命令在各场景下的行为}{Behavior of title line-break commands in different contexts}
    \label{tab:titlebreak}
    \begin{tabular}{@{}llll@{}}
        \toprule
        命令 & 封面标题 & 书脊标题 & 页眉右侧 \\ \midrule
        \verb|\titlebreak| & 换行 & 忽略 & 忽略 \\
        \verb|\\|           & 换行 & 忽略 & 换行 \\ \bottomrule
    \end{tabular}
\end{table}

使用时需注意以下两点：其一，页眉文字不得超过两行，本模板未对三行及以上页眉做适配，超出时排版效果将与 Word 不一致；其二，当标题在封面页因过长而发生自动换行时，自动换行产生的行间距偏窄（刻意为之），必须在适当位置手动添加换行符并指定行距。具体使用案例参见表~\ref{tab:titlebreak-examples}。

\begin{table}[H]
    \centering
    \bicaption{标题换行命令使用案例}{Usage examples of title line-break commands}
    \label{tab:titlebreak-examples}
    \zihao{-6}
    \begin{tabular}{@{}p{3cm}p{2.75cm}p{5cm}p{5cm}@{}}
        \toprule
        命令写法 & 封面标题 & 书脊（始终不换行） & 页眉 \\ \midrule
        \verb|\renewcommand{\cntitle}{...|\newline\verb|\titlebreak...}| &
        \parbox[t]{5cm}{大型企业与地区社会变迁研究\\——以广西为例} &
        大型企业与地区社会变迁研究——以广西为例 &
        大型企业与地区社会变迁研究——以广西为例 \\[2em]
        \verb|\renewcommand{\cntitle}{...|\newline\verb|\\...}| &
        \parbox[t]{5cm}{大型企业与地区社会变迁研究\\——以广西为例} &
        大型企业与地区社会变迁研究——以广西为例 &
        \parbox[t]{5cm}{大型企业与地区社会变迁研究\\——以广西为例} \\[2em]
        \verb|\renewcommand{\cntitle}{...|\newline\verb|\titlebreak...\\...}| &
        \parbox[t]{5cm}{大型企业\\与地区社会变迁研究\\——以广西为例} &
        大型企业与地区社会变迁研究——以广西为例 &
        \parbox[t]{5cm}{大型企业与地区社会变迁研究\\——以广西为例} \\
        \bottomrule
    \end{tabular}
\end{table}